\pagestyle{fancy}

\section{Methodology}
\label{sec:methodology}

This chapter defines the four-stage analytical pipeline used to test whether instruction tuning introduces new features into a language model's internal representations or predominantly reorganises pre-existing ones, mapping directly onto the research question of \S\ref{subsubsec:gap-rq}. The focus is methodological: what is measured, how the four-way classification is constructed, and which concrete settings determine the comparison. Stage one (Sections~\ref{subsec:models-saes}--\ref{subsec:activation-collection}) fixes the model pair, interpretability tool, and activation corpus, so that the pre-trained and instruction-tuned variants are compared at the same architecture, layer, and SAE width using matched Sparse Autoencoders. Stage two (Sections~\ref{subsec:feature-matching-classification}--\ref{subsec:threshold-selection}) defines a feature-level correspondence between the two checkpoints based on decoder direction similarity and activation correlation, and uses a threshold grid to check that the resulting four-way partition (preserved, modified, PT-exclusive, IT-exclusive) is not an artefact of a single threshold choice. Stage three (\S\ref{subsec:automated-labelling}) attaches descriptive vocabulary to sampled features through automated dual-LLM labelling with consensus construction, providing both the interpretive layer needed to discuss what features mean and an inter-annotator reliability check. Stage four (\S\ref{subsec:targeted-differential-analysis}) grounds the classification in behavioural prompts by computing targeted differential activations for a safety contrast (JailbreakBench harmful versus benign) and an instruction-following contrast (Stanford Alpaca versus TruthfulQA) using Welch's $t$-tests. The resulting artefacts are analysed in Chapters~\ref{sec:results} and \ref{sec:discussion}.

\subsection{Stage 1: Model Selection and Sparse Autoencoders}
\label{subsec:models-saes}

The experimental subject is the Gemma 3 1B model family from Google: the pre-trained base checkpoint \path{google/gemma-3-1b-pt} and its instruction-tuned counterpart \path{google/gemma-3-1b-it}. Both share the same architecture (26 transformer layers, hidden size 1152) and were loaded at \texttt{bfloat16} precision using HuggingFace \texttt{transformers}. Gemma 3 was chosen over larger alternatives for two reasons. First, the 1B scale fits both variants, their matched SAEs, and the activation buffers for a 5,000-sequence corpus into a single commodity GPU, a hard constraint for reproducibility and iteration speed. Second, this is the smallest scale at which publicly released SAEs exist for both a base and an instruction-tuned checkpoint of the same architecture, via Gemma Scope~2 \parencite{mcdougall_gemma_scope_2_2025}. Earlier project scoping targeted Gemma 2 2B, but that plan was abandoned because no matched SAE existed for the instruction-tuned checkpoint; without paired PT/IT SAEs at the same layer and width, the core comparison would not have been methodologically defensible.

The specific SAEs used here come from Gemma Scope 2, following prior SAE work that treats residual-stream activations as sparse combinations of interpretable feature directions \parencite{bricken_monosemanticity_2023,cunningham_sparse_2023,elhage_toy_2022}: \texttt{gemma-scope-2-1b-pt-res} for the base model and \texttt{gemma-scope-2-1b-it-res} for the instruction-tuned model, both at the \texttt{layer\_13\_width\_16k\_l0\_medium} configuration. Both have a dictionary width of 16,384 features, use the JumpReLU encoder of \textcite{rajamanoharan_improving_2024}, and were trained by Google DeepMind at comparable expected $L_0$ sparsity, making them directly comparable as feature dictionaries for the same architecture.

Layer 13 was chosen because it sits at the midpoint of the 26-layer stack, where prior work on Gemma Scope and residual-stream SAEs has reported the richest mixture of syntactic and semantic features. Middle layers are generally where abstract features relevant to downstream behaviour are most legible; earlier layers are dominated by token-level and syntactic features, later layers by output-oriented ones. This is a scoping decision, not a methodological claim, and its implications for generalisation are discussed in Chapter~\ref{sec:discussion}. One implementation detail worth recording: Gemma 3 1B required \texttt{bfloat16} rather than \texttt{float16} for inference, as the narrower dynamic range of \texttt{float16} produced overflow-driven NaN activations during early testing and would have silently corrupted the activation buffer.

\subsection{Stage 1: Activation Collection}
\label{subsec:activation-collection}

The activation corpus consists of 5,000 sequences drawn from the HuggingFace \path{HuggingFaceFW/fineweb} dataset (\path{sample-10BT} split), accessed in streaming mode with a fixed random seed for reproducibility. Each sequence was tokenised with the shared Gemma tokenizer and truncated to exactly 256 tokens; sequences shorter than 256 tokens were discarded rather than padded, so that every retained sequence contributes the same number of activation rows without requiring padding masks. The 5,000-sequence figure was downscaled from an initial 50,000-sequence plan after a pilot confirmed that the relevant SAE feature statistics (per-feature activation frequency, mean activation, per-sequence mean activation vectors used in \S\ref{subsec:feature-matching-classification}) stabilised well below 50,000 sequences, and that additional data mainly increased runtime without changing the qualitative results. The final corpus contains approximately 1.275 million activation rows per model (5,000 sequences $\times$ 255 tokens after the first token of each sequence is dropped, as described below).

For each of the two models, activations were collected from the residual stream output of layer 13 using a PyTorch forward hook registered on the corresponding \path{model.model.layers[13]} module. Sequences were processed in batches of four, and within each batch the first token of each sequence was dropped before SAE encoding, because the first residual-stream position of a causal transformer carries no context from other tokens and its activation statistics are atypical. The resulting activation matrix was moved to CPU memory in \texttt{float32} and then re-encoded through the SAE in chunks of 64 rows at a time; this chunked encoding was necessary because naively encoding the entire activation matrix in a single GPU call caused out-of-memory errors on the 16,384-wide feature dictionary, and the one-sequence-at-a-time alternative was too slow to be practical. For each chunk, the SAE encoder, decoder, and reconstruction residual were computed, and three running statistics were updated per feature: a sum of activations, a sum of squared activations, and a count of non-zero activations. These running statistics were used after all batches to compute the per-feature mean activation $\bar{f}_i$ and activation frequency $\rho_i$. In parallel, a per-sequence mean activation vector was computed for each sequence, producing a matrix of shape $5000 \times 16384$ that was later used as the input to the activation-correlation step of \S\ref{subsec:feature-matching-classification}. Finally, two SAE quality metrics were tracked throughout encoding: the fraction of variance unexplained (FVU) of the reconstruction and the mean $L_0$ count, computed as the average number of non-zero features per activation row. These were reported at the end of each encoding pass as sanity checks that the pre-trained SAEs were operating within their expected regime on the FineWeb corpus rather than failing silently.

The output of this phase, for each model, is a summary of SAE feature behaviour over the corpus: activation frequency, mean activation, per-sequence mean activation vectors, and the SAE decoder matrices (loaded once and reused). These four objects are the only inputs required for the feature matching and classification stage described next.

\subsection{Stage 2: Feature Matching and Four-Way Classification}
\label{subsec:feature-matching-classification}

The central methodological construct is the four-way classification of SAE features into preserved, modified, PT-exclusive, and IT-exclusive categories, built from two complementary similarity measures: decoder direction similarity and activation correlation. Decoder similarity captures whether the two dictionaries point in similar directions in residual-stream space, while activation correlation captures whether matched features are used similarly on the shared FineWeb corpus. Using both measures jointly, and requiring agreement between them for the strongest category (preserved), reduces the weaknesses of either measure in isolation.

Decoder direction similarity is computed from the PT and IT decoder matrices, $D^{\mathrm{PT}} \in \mathbb{R}^{F \times d}$ and $D^{\mathrm{IT}} \in \mathbb{R}^{F \times d}$, after row-wise $\ell_2$ normalisation. The cosine similarity matrix is then $S = D^{\mathrm{PT}} (D^{\mathrm{IT}})^{\top}$. For each PT feature $i$, the best matching IT feature is $j^{*}(i) = \arg\max_j S_{ij}$, with maximum cosine similarity $s^{\mathrm{max}}_i$; the same procedure is run in reverse for IT features. This bidirectional matching matters because PT-to-IT and IT-to-PT nearest neighbours need not coincide.

Activation correlation is computed over the per-sequence mean activation matrices collected in \S\ref{subsec:activation-collection}. For each feature and its best geometric match, the Pearson correlation of their per-sequence activation vectors is computed over the 5,000 sequences, with a numerical guard for zero-variance features. The resulting scalar $\rho_i \in [-1,1]$ measures whether a feature and its nearest neighbour in the other dictionary are functionally correlated on the same inputs; the same quantity is computed in the reverse direction for IT features.

With these two measures in hand, the four-way classification at thresholds $(\tau_d,\tau_a)$ is defined as follows. A PT feature $i$ is \textit{preserved} if both $s^{\mathrm{max}}_i \geq \tau_d$ and $\rho_i \geq \tau_a$; \textit{modified} if exactly one condition holds; and \textit{PT-exclusive} if neither holds. The fourth category, \textit{IT-exclusive}, is defined symmetrically on IT features using the reverse-direction matches and their activation correlations. Analysis is restricted to \textit{active} features with activation frequency above 0.1\% of tokens, which removes dead features without materially constraining the category distribution.

\subsection{Stage 2: Threshold Selection and Sensitivity Analysis}
\label{subsec:threshold-selection}

The classification depends on two free parameters, $\tau_d$ and $\tau_a$, so the full procedure is recomputed on a 25-point grid with $\tau_d \in \{0.5,0.6,0.7,0.8,0.9\}$ and $\tau_a \in \{0.1,0.2,0.3,0.5,0.7\}$. For each pair, the preserved, modified, PT-exclusive, and IT-exclusive counts are recorded together with preservation and modification rates, and summarised as heatmaps. The aim is robustness rather than optimisation: to show that the qualitative shape of the classification is stable across reasonable threshold choices.

The primary operating point for the interpretive phases (Subsections~\ref{subsec:automated-labelling} and \ref{subsec:targeted-differential-analysis}) is $\tau_d=0.7$, $\tau_a=0.3$, selected after the sensitivity grid was computed because it lies in the interior of the grid and represents a moderate standard of both geometric and functional correspondence. Stricter settings collapse the preserved set and inflate the exclusive categories, while more permissive settings blur the distinction between modified and preserved. The full grid in \S\ref{subsec:feature-classification-threshold-robustness} serves as the robustness check.

\subsection{Stage 3: Automated Feature Labelling and Consensus Construction}
\label{subsec:automated-labelling}

A four-way classification without descriptions of what the features actually detect would be hard to interpret. To attach a descriptive vocabulary while keeping the effort manageable, this stage samples 50 features from each category, for 200 labelled features in total, and submits their top-activating contexts to two independent LLMs, Claude Haiku 4.5 \parencite{anthropic_haiku_2025} and GPT-5.4-nano, for automated labelling. Using two models rather than one is deliberate: a single model's labels are difficult to audit, while two independent models provide an inter-annotator reliability check and make labeller disagreement an observable that can itself be analysed.

Within each category, features are sampled with probability proportional to their activation frequency rather than uniformly. This choice weights the sample toward features that contribute meaningfully to the model's computation over the corpus, and avoids allocating labelling budget to rarely-active features whose top-activating examples are noisy and hard to interpret. For each sampled feature, the top 20 activating token contexts are collected by running the appropriate model and SAE over the first 1,000 sequences of the 5,000-sequence FineWeb corpus used in Stage 1, extracting per-token feature activations, and keeping the highest-activating positions together with a window of 8 tokens on either side of the activating token. The activating token itself is marked with \texttt{>>>} and \texttt{<<<} delimiters in the extracted context so that the labelling model can unambiguously identify which token the feature fires on. This delimiter convention follows the display format used in the original Anthropic monosemanticity work and in subsequent SAE labelling pipelines.

Both models receive identical prompts and identical output schemas, so that any differences in output can be attributed to the models rather than prompt variation. The prompt instructs the model to act as a neural network interpretability expert, shows the 20 top-activating contexts for a feature, and asks for three outputs: a one-to-two sentence description of the detected pattern, a single-label category from a six-way taxonomy (\path{language_syntax}, \path{knowledge_semantics}, \path{instruction_style}, \path{safety_alignment}, \path{formatting}, \path{other}), and a confidence rating from 1 to 5. Output is constrained to this schema using the structured-output interfaces of the Anthropic and OpenAI SDKs (\path{client.messages.parse} with a Pydantic \texttt{FeatureLabel} schema for Haiku, \path{client.beta.chat.completions.parse} with the equivalent schema for GPT). Structured outputs were essential: an earlier free-text version produced inconsistent category strings (``syntactic'', ``syntax'', ``grammar'', ``language'') requiring ad-hoc normalisation, while the schema-constrained version produces categories directly comparable across models. Labelling ran category-by-category with a rate-limit sleep between calls, and labels were serialised to disk as JSON after each category completed.

Consensus between the two labelling models is constructed as follows. For each sampled feature, if both models assign the same category, the feature is marked as an agreement case and that category becomes its consensus label. If the two models assign different categories, the feature is marked as a disagreement case, and the consensus label is taken from whichever model reported the higher confidence rating for that feature; in the event of equal confidence, Haiku is used as the primary model and GPT as the secondary. The consensus record for each feature additionally retains both descriptions and both confidence ratings, so that disagreement cases can be inspected qualitatively rather than only counted. The resulting consensus labels support two distinct uses in Chapter~\ref{sec:results}: first, the agreement rate and disagreement confusion matrix serve as an inter-annotator reliability diagnostic and give a measurable lower bound on the stability of the six-way taxonomy; second, within the agreement set, the distribution of categories across the four feature classes provides the descriptive evidence used to discuss what kinds of features are preserved, modified, or exclusive to each checkpoint. Both uses are discussed as results in \S\ref{subsec:automated-labelling-inter-annotator-agreement}.

\subsection{Stage 4: Targeted Differential Analysis on Behavioural Probes}
\label{subsec:targeted-differential-analysis}

The labelling stage attaches descriptive vocabulary to features but does not by itself show that any of them matters for the behaviours instruction tuning is supposed to induce. To connect the classification to behaviour, a targeted differential analysis is run using two behavioural contrasts instantiated with three external probe datasets: JailbreakBench \parencite{chao_jailbreakbench_2024} for the safety contrast, Stanford Alpaca \parencite{taori_stanford_alpaca_2023} for instruction-shaped prompts, and TruthfulQA \parencite{lin_truthfulqa_2022} as the factual contrast set for the instruction-following comparison. In the concrete study reported here, the safety contrast uses 100 harmful and 100 benign JailbreakBench prompts, and the instruction contrast uses the first 200 Alpaca instructions against 200 TruthfulQA questions from the generation split. The contrast design matters. For safety, harmful-split JailbreakBench prompts are contrasted against benign-split prompts from the same dataset rather than against generic FineWeb text, so that the differential activation isolates safety-relevant structure rather than incidental stylistic differences between task-oriented prompts and web text. For instructions, the Alpaca prompts are contrasted against TruthfulQA questions to produce a contrast between imperatively phrased task prompts and interrogatively phrased factual queries.

For each contrast and each of the two models, SAE feature activations are computed on every probe prompt using the same activation collection procedure as in \S\ref{subsec:activation-collection}, but with a shorter maximum sequence length of 128 tokens appropriate to the probe prompts and with per-prompt rather than per-sequence feature statistics. The per-prompt activation of a feature is defined as the mean of its token-level activations over all non-first tokens of the prompt, and is stored as a row in a prompts $\times$ features matrix per model. Two such matrices are produced per model---one for the safety-vs-benign contrast and one for the instruction-vs-truthful contrast---and all subsequent analysis is performed feature-by-feature on these matrices.

For each feature, a Welch's two-sample $t$-test \parencite{welch_generalization_1947} is run comparing its per-prompt activations on the target group (e.g.\ harmful prompts) to those on the contrast group (e.g.\ benign prompts). Welch's variant is used rather than the standard equal-variance $t$-test because the two groups cannot be assumed to have equal variance---a feature that is selectively active on harmful prompts is, almost by definition, more variable on the group where it activates than on the group where it does not. The implementation uses \texttt{scipy.stats.ttest\_ind} with \texttt{equal\_var=False}. Features are retained as differential if they satisfy two conditions simultaneously: a $t$-test $p$-value below 0.01 and a target-group mean activation above 0.01, the latter being a noise-floor filter that prevents the analysis from being dominated by features that are statistically elevated only because their baseline activation is essentially zero. For each retained feature the pipeline records its feature index, target-group and contrast-group mean activations, their difference and ratio, and the test statistic and $p$-value. The top 30 differential features per contrast per model are then ranked by difference and carried forward to the labelling pipeline of \S\ref{subsec:automated-labelling}, which re-labels them on examples drawn from the probe prompts themselves rather than from FineWeb. This re-labelling step produces descriptions that are grounded in the behavioural contrast rather than in generic web text, and enables the results of \S\ref{subsec:targeted-differential-analysis-results} to discuss not only which features are differential but what they appear to detect.
